%%% Основные сведения %%%
\newcommand{\thesisAuthorLastName}{\fixme{Фамилия}}
\newcommand{\thesisAuthorOtherNames}{\fixme{Имя Отчество}}
\newcommand{\thesisAuthorInitials}{\fixme{И.\,О.}}
\newcommand{\thesisAuthor}             % Диссертация, ФИО автора
{%
    \texorpdfstring{% \texorpdfstring takes two arguments and uses the first for (La)TeX and the second for pdf
        \thesisAuthorLastName~\thesisAuthorOtherNames% так будет отображаться на титульном листе или в тексте, где будет использоваться переменная
    }{%
        \thesisAuthorLastName, \thesisAuthorOtherNames% эта запись для свойств pdf-файла. В таком виде, если pdf будет обработан программами для сбора библиографических сведений, будет правильно представлена фамилия.
    }
}
\newcommand{\thesisAuthorShort}        % Диссертация, ФИО автора инициалами
{\thesisAuthorInitials~\thesisAuthorLastName}
%\newcommand{\thesisUdk}                % Диссертация, УДК
%{\fixme{xxx.xxx}}
\newcommand{\thesisTitle}              % Диссертация, название
{\fixme{Длинное название диссертационной работы, состоящее из~достаточно большого
количества слов, совсем длинное длинное длинное длинное название, из~которого
простому обывателю знакомы, в~лучшем случае, лишь отдельные слова}}
\newcommand{\thesisSpecialtyNumber}    % Диссертация, специальность, номер
{\fixme{XX.XX.XX}}
\newcommand{\thesisSpecialtyTitle}     % Диссертация, специальность, название (название взято с сайта ВАК для примера)
{\fixme{Технология обработки, хранения и~переработки злаковых, бобовых культур,
крупяных продуктов, плодоовощной продукции и~виноградарства}}
%% \newcommand{\thesisSpecialtyTwoNumber} % Диссертация, вторая специальность, номер
%% {\fixme{XX.XX.XX}}
%% \newcommand{\thesisSpecialtyTwoTitle}  % Диссертация, вторая специальность, название
%% {\fixme{Теория и~методика физического воспитания, спортивной тренировки,
%% оздоровительной и~адаптивной физической культуры}}
\newcommand{\thesisDegree}             % Диссертация, ученая степень
{\fixme{кандидата физико-математических наук}}
\newcommand{\thesisDegreeShort}        % Диссертация, ученая степень, краткая запись
{\fixme{канд. физ.-мат. наук}}
\newcommand{\thesisCity}               % Диссертация, город написания диссертации
{\fixme{Город}}
\newcommand{\thesisYear}               % Диссертация, год написания диссертации
{\the\year}
\newcommand{\thesisOrganization}       % Диссертация, организация
{\fixme{Федеральное государственное автономное образовательное учреждение высшего
образования <<Длинное название образовательного учреждения <<АББРЕВИАТУРА>>}}
\newcommand{\thesisOrganizationShort}  % Диссертация, краткое название организации для доклада
{\fixme{НазУчДисРаб}}

\newcommand{\thesisInOrganization}     % Диссертация, организация в предложном падеже: Работа выполнена в ...
{\fixme{учреждении с~длинным длинным длинным длинным названием, в~котором
выполнялась данная диссертационная работа}}

%% \newcommand{\supervisorDead}{}           % Рисовать рамку вокруг фамилии
\newcommand{\supervisorFio}              % Научный руководитель, ФИО
{\fixme{Фамилия Имя Отчество}}
\newcommand{\supervisorRegalia}          % Научный руководитель, регалии
{\fixme{уч. степень, уч. звание}}
\newcommand{\supervisorFioShort}         % Научный руководитель, ФИО
{\fixme{И.\,О.~Фамилия}}
\newcommand{\supervisorRegaliaShort}     % Научный руководитель, регалии
{\fixme{уч.~ст.,~уч.~зв.}}

%% \newcommand{\supervisorTwoDead}{}        % Рисовать рамку вокруг фамилии
%% \newcommand{\supervisorTwoFio}           % Второй научный руководитель, ФИО
%% {\fixme{Фамилия Имя Отчество}}
%% \newcommand{\supervisorTwoRegalia}       % Второй научный руководитель, регалии
%% {\fixme{уч. степень, уч. звание}}
%% \newcommand{\supervisorTwoFioShort}      % Второй научный руководитель, ФИО
%% {\fixme{И.\,О.~Фамилия}}
%% \newcommand{\supervisorTwoRegaliaShort}  % Второй научный руководитель, регалии
%% {\fixme{уч.~ст.,~уч.~зв.}}

\newcommand{\opponentOneFio}           % Оппонент 1, ФИО
{\fixme{Фамилия Имя Отчество}}
\newcommand{\opponentOneRegalia}       % Оппонент 1, регалии
{\fixme{доктор физико-математических наук, профессор}}
\newcommand{\opponentOneJobPlace}      % Оппонент 1, место работы
{\fixme{Не очень длинное название для места работы}}
\newcommand{\opponentOneJobPost}       % Оппонент 1, должность
{\fixme{старший научный сотрудник}}

\newcommand{\opponentTwoFio}           % Оппонент 2, ФИО
{\fixme{Фамилия Имя Отчество}}
\newcommand{\opponentTwoRegalia}       % Оппонент 2, регалии
{\fixme{кандидат физико-математических наук}}
\newcommand{\opponentTwoJobPlace}      % Оппонент 2, место работы
{\fixme{Основное место работы c длинным длинным длинным длинным названием}}
\newcommand{\opponentTwoJobPost}       % Оппонент 2, должность
{\fixme{старший научный сотрудник}}

%% \newcommand{\opponentThreeFio}         % Оппонент 3, ФИО
%% {\fixme{Фамилия Имя Отчество}}
%% \newcommand{\opponentThreeRegalia}     % Оппонент 3, регалии
%% {\fixme{кандидат физико-математических наук}}
%% \newcommand{\opponentThreeJobPlace}    % Оппонент 3, место работы
%% {\fixme{Основное место работы c длинным длинным длинным длинным названием}}
%% \newcommand{\opponentThreeJobPost}     % Оппонент 3, должность
%% {\fixme{старший научный сотрудник}}

\newcommand{\leadingOrganizationTitle} % Ведущая организация, дополнительные строки. Удалить, чтобы не отображать в автореферате
{\fixme{Федеральное государственное бюджетное образовательное учреждение высшего
профессионального образования с~длинным длинным длинным длинным названием}}

\newcommand{\defenseDate}              % Защита, дата
{\fixme{DD mmmmmmmm YYYY~г.~в~XX часов}}
\newcommand{\defenseCouncilNumber}     % Защита, номер диссертационного совета
{\fixme{Д\,123.456.78}}
\newcommand{\defenseCouncilTitle}      % Защита, учреждение диссертационного совета
{\fixme{Название учреждения}}
\newcommand{\defenseCouncilAddress}    % Защита, адрес учреждение диссертационного совета
{\fixme{Адрес}}
\newcommand{\defenseCouncilPhone}      % Телефон для справок
{\fixme{+7~(0000)~00-00-00}}

\newcommand{\defenseSecretaryFio}      % Секретарь диссертационного совета, ФИО
{\fixme{Фамилия Имя Отчество}}
\newcommand{\defenseSecretaryRegalia}  % Секретарь диссертационного совета, регалии
{\fixme{д-р~физ.-мат. наук}}            % Для сокращений есть ГОСТы, например: ГОСТ Р 7.0.12-2011 + http://base.garant.ru/179724/#block_30000

\newcommand{\synopsisLibrary}          % Автореферат, название библиотеки
{\fixme{Название библиотеки}}
\newcommand{\synopsisDate}             % Автореферат, дата рассылки
{\fixme{DD mmmmmmmm}\the\year~года}

% To avoid conflict with beamer class use \providecommand
\providecommand{\keywords}%            % Ключевые слова для метаданных PDF диссертации и автореферата
{}

%%% Данные выпускной квалификационной работы %%%
\newcommand{\vkrMinistry}{Министерство науки и высшего образования Российской Федерации}
\newcommand{\vkrUniversityType}{Федеральное государственное автономное образовательное учреждение}
\newcommand{\vkrUniversityLevel}{высшего образования}
\newcommand{\vkrUniversityName}{<<Волгоградский государственный университет>>}
\newcommand{\vkrUniversity}{\vkrUniversityType\ \vkrUniversityLevel\ \vkrUniversityName}
\newcommand{\vkrInstitute}{институт Математики и информационных технологий}
\newcommand{\vkrDepartment}{кафедра Фундаментальной информатики и искусственного интеллекта}
\newcommand{\vkrDepartmentShort}{ФИИИ}
\newcommand{\vkrDepartmentHeadTitle}{Зав. каф. \vkrDepartmentShort}
\newcommand{\vkrDepartmentHeadDegree}{д.ф.-м.н., профессор}
\newcommand{\vkrDepartmentHead}{Воронин Александр Александрович}
\newcommand{\vkrDepartmentHeadShort}{А.\,А.~Воронин}
\newcommand{\vkrStudentFullName}{Савкин Никита Дмитриевич}
\newcommand{\vkrStudentGroup}{ПМИб-221}
\newcommand{\vkrWorkType}{бакалаврская работа}
\newcommand{\vkrWorkKind}{ВЫПУСКНАЯ КВАЛИФИКАЦИОННАЯ РАБОТА}
\newcommand{\vkrDirectionCode}{01.03.02}
\newcommand{\vkrDirectionTitle}{Прикладная математика и информатика}
\newcommand{\vkrProfile}{Создание и применение технологий больших данных}
\newcommand{\vkrTitle}{Разработка интеллектуальной системы для классификации текстовых документов}
\newcommand{\vkrSupervisorDegree}{Кандидат технических наук}
\newcommand{\vkrSupervisorPosition}{доцент кафедры ФИИИ}
\newcommand{\vkrSupervisorFullName}{Васильченко Анна Анатольевна}
\newcommand{\vkrSupervisorShort}{А.\,А.~Васильченко}
\newcommand{\vkrCity}{Волгоград}
\newcommand{\vkrYear}{\the\year}
\newcommand{\vkrDefenseDay}{\fixme{ДД}}
\newcommand{\vkrDefenseMonth}{\fixme{месяца}}
\newcommand{\vkrDefenseProtocol}{\fixme{номер}}
\newcommand{\vkrAssignmentApprovedDay}{\fixme{ДД}}
\newcommand{\vkrAssignmentApprovedMonth}{\fixme{месяца}}
\newcommand{\vkrAssignmentApprovedYear}{\the\year}
\newcommand{\vkrAssignmentIssueDate}{\fixme{ДД.ММ.ГГГГ}}
\newcommand{\vkrAssignmentDeadline}{\fixme{ДД.ММ.ГГГГ}}
\newcommand{\vkrAssignmentGoal}{исследовать и~разработать локальную систему
классификации электронной почты методами обработки естественного языка,
векторизации, кластеризации и~машинного обучения}
\newcommand{\vkrAssignmentTasks}{%
    \item проанализировать методы обработки текстов и~машинного обучения
        для~классификации электронной почты;
    \item разработать и~реализовать процесс предобработки, векторизации,
        кластеризации и~классификации сообщений;
    \item сформулировать постановку задачи и~оценить качество системы
        по~метрикам классификации, покрытию и~доле отказов.
}
\newcommand{\vkrAssignmentStages}{%
    \item обзор литературы, постановка задачи и~выбор методов;
    \item проектирование, реализация и~экспериментальная оценка системы;
    \item оформление ВКР и~подготовка к~защите.
}
\newcommand{\vkrAssignmentLiterature}{%
    \item Большакова~Е.\,И., Клышинский~Э.\,С., Ландэ~Д.\,В.,
        Носков~А.\,А., Пескова~О.\,В., Ягунова~Е.\,В. Автоматическая
        обработка текстов на~естественном языке и~компьютерная лингвистика:
        учебное пособие. --- М.: МИЭМ, 2011. --- 272~с.
    \item Лукашевич~Н.\,В., Сорокин~А.\,А. Компьютерная лингвистика
        и~автоматическая обработка текстов: учебник. --- М.: МАКС Пресс,
        2025. --- 601~с.
    \item Батура~Т.\,В. Математическая лингвистика и~автоматическая обработка
        текстов на~естественном языке: учебное пособие. --- Новосибирск:
        РИЦ НГУ, 2016. --- 166~с.
    \item Журавлёв~Ю.\,И., Рязанов~В.\,В., Сенько~О.\,В. Распознавание:
        математические методы, программная система, практические применения.
        --- М.: Фазис, 2006. --- 176~с.
    \item Sebastiani~F. Machine Learning in Automated Text Categorization //
        ACM Computing Surveys. --- 2002. --- Vol.~34, no.~1. --- P.~1--47.
}

\providecommand{\documentPdfTitle}{\thesisTitle}
\providecommand{\documentPdfAuthor}{\thesisAuthor}
\providecommand{\documentPdfSubject}{\thesisSpecialtyNumber\ \thesisSpecialtyTitle}
\renewcommand{\documentPdfTitle}{\vkrTitle}
\renewcommand{\documentPdfAuthor}{\vkrStudentFullName}
\renewcommand{\documentPdfSubject}{\vkrDirectionCode\ \vkrDirectionTitle}
